\renewcommand{\arraystretch}{1.3}
\captionsetup[table]{justification=raggedright,singlelinecheck=false}
\begin{table}[H]
\caption{Pemetaan \textit{Intent Type} terhadap Kedalaman Traversal Graf}
\label{tbl:intent_depth_iv}
\centering
\begin{tabular}{|p{0.22\textwidth}|p{0.10\textwidth}|p{0.55\textwidth}|}
\hline
\textbf{\textit{Intent Type}} & \textbf{Depth} & \textbf{Alasan} \\
\hline
\texttt{explain} & 2 & Cukup untuk menjawab definisi dan cara kerja \textit{resource}; depth 3+ menambahkan \textit{noise} generik. \\
\hline
\texttt{followup} & 2 & Pertanyaan lanjutan biasanya merujuk entitas yang sama; konteks depth 2 sudah tersedia dari giliran sebelumnya. \\
\hline
\texttt{generate\_yaml} & 3 & Pembuatan YAML membutuhkan field tingkat kontainer (\texttt{image}, \texttt{ports}, \texttt{env}) yang berada di depth 3. \\
\hline
\texttt{trace\_relationship} & 3 & Penelusuran relasi antar-\textit{resource} menjangkau jembatan lintas-\textit{resource} pada depth 2--3. \\
\hline
\end{tabular}
\end{table}
